\section{教育背景}

\textbf{电子科技大学 (UESTC)} \hfill 2022年9月 --- 2026年6月

\textbf{格拉斯哥大学 (UofG), 双学位项目} \hfill 2022年9月 --- 2026年6月

\begin{itemize}
    \item \textbf{专业学位}: 电子与信息工程 (ECE) 学士; \MYhref{https://marcobisky.github.io/cv/score.pdf}{GPA: 3.87/4.0}, \MYhref{https://marcobisky.github.io/cv/rank.pdf}{排名: 2/164 (前1.2\%)}.
    \item \textbf{相关课程}: 信息论, 随机过程, 流匹配与扩散模型, 大语言模型中的强化学习等.
\end{itemize}

\textbf{香港中文大学 (深圳)} \hfill 2026年9月

\begin{itemize}
    \item \textbf{专业学位}: 计算机科学与技术 (CS) 研究型硕士
\end{itemize}


\section{研究与项目经历}

\MYhref{https://github.com/Marcobisky/llmlab}{\textbf{LLMlab: 基于形式语言的 LLM RLVR 快速后训练验证平台}} \hfill 2026年6月

\begin{itemize}
    \item 提出动机: 为了降低 LLM 的极高算法验证成本, LLMlab 可在可控的合成数据上进行高效的算法性能验证.
    \item 设计了一种确定性的形式语言, 支持在不同形式语言难度上的 RLVR 的高效算法性能验证.
    \item 实现了在 2.67M、0.15M 的教师和学生模型上的完整 pretrain, SFT, GRPO, KD, OPD, SDPO 流水线和平行消融实验, 最高实现了在语言难度 0-3 等级的 100\% 正确率和难度 4-6 等级超过 70\% 的正确率.
    \item 集成了 PCA 投影的 loss landscape 及权重变化轨迹、逐层注意力热力图、exposure bias 度量等可视化工具集以便对模型的可解释性和训练过程进行深入分析.
\end{itemize}

\textbf{LLM 后训练: 复现 \MYhref{https://arxiv.org/abs/2601.20802}{Self-Distillation Policy Optimization (SDPO)} 论文} \hfill 2026年5月

\begin{itemize}
    \item 在理解 SDPO 算法的基础上, 基于 verl RL 框架在 RunPod H100 GPU 上成功复现 SDPO 算法.
    \item 使用 Qwen2.5-3B 在代码生成任务上训练, 利用 LeetCode 风格的反馈 (运行时错误、失败测试用例) 作为学习信号, 通过 WandB 跟踪 40 step 的 SDPO 训练动态, 记录平均 reward、相对参考 policy 的 token 级 KL 散度.
\end{itemize}

\textbf{边缘 TinyML 的 RISCV 加速器系统级软硬件协同设计} \hfill 2025年9月 --- 2026年4月

\vspace{-0.2em} \begin{small} \textit{研究助理, \MYhref{https://scholar.google.com/citations?user=1NT1jFMAAAAJ&hl=en}{李耘 教授}, 电子科技大学} \end{small} \vspace{-0.2em}

\begin{itemize}
    \item 在 Artix-7 FPGA 上独立架构并开发了用于实时 YOLOv8n 边缘推理的神经网络处理单元 (NPU), 通过自定义 RISC-V 指令集扩展 (Xnpu) 实现了无软核处理器的纯硬件执行架构.
    \item 自研端到端 Python 机器学习编译器，实现了全自动 INT16 量化与内存感知的指令调度，将量化后的 mAP 精度损失严格控制在 PyTorch FP32 基线的 0.3\% 以内.
    \item 设计了基于 3×3 脉动阵列 (Systolic Array) 的参数化 RTL 算子，并实现了全硬件加速的后处理引擎 (DFL 框解码与 NMS)，峰值算力达 288 MACs/cycle 与 23.4 GOPS.
    \item 构建了完整的边缘计算端到端数据流，集成了异步摄像头/UDP 视频流管道及 AXI4 DDR3L 内存总线复用，并基于 \texttt{Cocotb}、\texttt{Bazel} 和 \texttt{Icarus Verilog} 完成了全系统的仿真与验证.
\end{itemize}

\MYhref{https://github.com/Marcobisky/my-riscv}{\textbf{完整单周期 RV32I CPU 内核的设计与可视化}} \hfill 2025年1月 --- 2025年3月

\begin{itemize}
    \item 使用 \texttt{Verilog} 从零手搓了一个单核、单周期 32 位 RISCV CPU, 使用 \texttt{Verilator} 仿真, 并在 \texttt{Digital} 软件中进行了工作原理可视化, 已在 \MYhref{https://github.com/Marcobisky/my-riscv}{Github} 上开源.
    \item 构建了完整的数据通路, 包括 PC、取指器、译码器、寄存器堆、ALU、基于 LRU 的 L1 缓存等, 支持 GPIO, IIC, UART 等基本外设.
    \item 使用 RISCV 汇编实现了一个引导程序, 以及使用 \texttt{C} 语言实现了基本的 delay 和 GPIO 库. 使用 RISCV GNU 工具链进行交叉编译和仿真.
\end{itemize}


\section{相关技能}

\begin{tabularx}{\linewidth}{@{}l X@{}}
\textbf{编程语言} &  \normalsize{\texttt{Python}, \texttt{PyTorch}, \texttt{C/C++} , \texttt{Makefile}.}\\
\textbf{英语} &  \normalsize{英语四六级 (615, 573), \MYhref{https://marcobisky.github.io/cv/ielts-score.pdf}{雅思 7} (8/7/7/6.5)}\\
\end{tabularx}


\section{奖项}

\textbf{电子科技大学一等学业奖学金 (前 5\%)} \hfill 2023年12月, 2024年12月 \\
\textbf{国家奖学金 (前 0.2\%)} \hfill 2024年12月 \\
\textbf{\MYhref{https://www.gla.uestc.edu.cn/info/1089/17553.htm}{2026 届省级优秀毕业生}} \hfill 2025年10月

